\section{Discussion}
\label{sec:discussion}
\subsection{Separating auxiliary supervision from forecast readout}
The mean of the five separately evaluated seed-44 fold models gives a descriptive hierarchy in Macro AUROC and AP: direct future-state supervision improves the direct-readout means, and forecast-attention readout increases them further. Direct readout nevertheless has the highest mean Macro F1. Forecast attention does not lead AUROC at T0 and has lower Macro AUROC than both direct controls in fold 5. None of the nine principal Macro comparisons meets the adjusted criterion. The results therefore describe the selected fitted models without establishing architectural superiority or equivalence.

The supplementary mean-pooling comparison also remains uncertain. Attention has higher mean Macro AUROC, AP, and F1, but their paired intervals include zero, and the landmark ordering is not uniform. This does not establish that attention is necessary. The \direct{}--\nojepa{} contrast most closely isolates future-state supervision; the full-versus-direct contrast changes forecast features, pooling, and classifier width together. Classifier-width-matched direct readout, stopped-gradient forecasts, copy-state readout, and independent holdout forecast evaluation would help distinguish these contributions.

The main models predict embeddings rather than generate MRI for classification. The internal name \nojepa{} denotes removal of the future-state objective while retaining the shared fusion network and transition queries. Archived epoch-20 probes show higher copy-normalized forecast skill for auxiliary-supervised direct models than for forecast attention at every horizon, despite the latter's higher mean pCR AUROC. A possible mechanism is that endpoint gradients shape forecasts toward classification-relevant directions at the expense of some EMA-target fidelity. Gradient routing is consistent with that possibility, but training probes in separately learned spaces do not establish it or demonstrate anatomical forecast quality.

\subsection{What additional visits change}
Forecast attention's mean AUROC increases across T0--T2, with a positive unadjusted T2--T0 interval. The corresponding test does not pass secondary-family or all-90 correction, and the AP and F1 intervals include zero. The T1 operating point increases sensitivity while adding false positives and reducing specificity. Added visits, ranking, and a fixed-threshold decision rule therefore need not improve together. Outcome-conditioned histograms summarize within-fold probability changes and do not establish improved predictions for every patient or individual treatment effects.

T0 imaging preceded randomization, while the model is conditioned on the eventual assigned arm. This remains a retrospective conditional scenario. The arm encoding does not capture delivered dose, cycle timing, dose reductions, adherence, or treatment changes. Altering the arm token cannot identify an individual treatment effect without suitable causal assumptions and exposure data. Complete T0--T2 holdout eligibility also restricts the population to which the longitudinal comparison applies.

\subsection{Classification, probability quality, and hypothetical utility}
The additional metrics reveal operating trade-offs that are obscured by AUROC. Forecast attention has higher mean sensitivity and NPV but lower specificity and PPV than auxiliary-supervised direct readout. Relative to \nojepa{}, both the sensitivity increase and specificity decrease meet the secondary-family and all-90 adjusted criteria. Its specificity decrease relative to direct readout also meets both corrections. These are conditional findings about fixed fitted models at threshold 0.5, and the adverse specificity result should accompany the sensitivity gain. All principal models' Macro raw accuracies are below the all-negative reference, and direct readout has the highest mean balanced accuracy and MCC among them.

Forecast attention also has worse NLL, Brier score, and ECE point estimates than both direct controls. Discrimination ordering therefore does not establish probability calibration. Selecting a Youden threshold from each matching fold validation set increases mean F1 only for \nojepa{} and decreases it for the other models. These thresholds avoid combining incompatible score definitions across models, but the validation patients also informed checkpoint selection, preventing an independent or nested calibration claim.

Net-benefit orderings change with the hypothetical threshold. Direct readout has higher mean Macro net benefit at 0.30 and 0.50, while forecast attention is higher at 0.70; every principal variant has negative mean Macro net benefit at 0.50. These are fold means under a generic positive-action utility, not evaluations of a defined care pathway. A clinical threshold requires the consequences of false-positive and false-negative actions to be specified and calibration to be assessed without holdout leakage.

\subsection{Fold variation and subgroup interpretation}
The mean and sample standard deviation across folds describe five models fitted to overlapping development partitions under one seed. Because these models reuse the same holdout, fold standard deviations do not supply independent clinical replications, training-seed variability, or patient standard errors. Shared-patient-bootstrap inference instead measures conditional patient-sampling uncertainty for the average metric of these fixed models. Folds are not resampled and models are not refitted, so those intervals do not characterize the whole training procedure.

HR/HER2 subgroup intervals all include zero; the HR-positive/HER2-positive point estimate favors \nojepa{}, whereas the other three groups favor forecast attention descriptively. No subgroup-interaction hypothesis is tested, and the smallest group contains only 14 patients. Source namespaces share one clinical release and do not supply external validation. The crop-clipping subgroup contains 11 patients and one positive outcome; suppressing its AUROC/AP is more informative than interpreting an unstable discrimination estimate.

\subsection{Limitations and a confirmatory study}
\label{sec:limitations}
Seed 44, separate-fold metric averaging, and the reporting procedures were selected after holdout results were available. Patient-bootstrap inference conditions on these fitted models and fixed validation-selected thresholds. It excludes retraining, seed choice, model selection, aggregation choice, and threshold-selection uncertainty. One selected seed does not establish training stability, and five overlapping training folds are not a substitute for independent training-seed assessment. The reported mean is the average performance of five distinct predictors, not a deployed combined predictor or a model retrained on the full development cohort.

The 95\% patient intervals are unadjusted. Holm corrections within nine principal and 81 secondary comparisons, together with the all-90 sensitivity analysis, do not undo retrospective choices. Null-centered bootstrap tests are approximate, particularly for non-smooth statistics such as AP, F1, and binned ECE. The two specificity tests reached the Monte Carlo resolution floor, which limits precision for their raw $p$ values. Nonsignificant discrimination contrasts do not establish equivalence. A confirmatory study should lock training seeds, aggregation, model and threshold selection, and statistical comparisons before evaluating independent patients.

The holdout is an internal downstream split. Twenty-nine patients had belonged to an earlier development split. Current supervised models were retrained using the new patient-disjoint folds and were not warm-started from those earlier pCR models, but codec and localization-model pretraining have not been fully audited for current-holdout overlap. This prevents a claim that every pipeline stage was evaluated on wholly unseen patients.

Fixed baseline localization and cropping may omit multifocal disease or changes outside the field of view. Geometry and array checks do not establish anatomical registration accuracy. Holdout patients have complete T0--T2 imaging, whereas development retains incomplete histories; missingness, acquisition variation, and subtype prevalence may limit transportability. Archived fixed-training-probe diagnostics do not replace best-checkpoint holdout forecast evaluation and have small long-horizon denominators.

The study remains an internal architectural comparison. Matched clinical-only, conventional imaging, and independent longitudinal baselines are still needed, along with external validation, anatomical quality review, and independently assessed calibration. Precise pCR pathology criteria, scan intervals, ethics determination, and upstream provenance also require final documentation before submission.
